\section{Fixed certificate for the two initial batches}
\label{sec:two-phase-certificate}

The local proof uses a comparison distribution to show that some
candidate schedule has enough combined first-hit and mode-saving
credit while preserving the all-negative waiting bound. This
distribution certifies existence; the algorithm searches candidate
schedules and minimizes their rational score. We list the complete
fixed distribution here, making the calibration reproducible directly
from the stated data.

This table specifies the proof distribution in
Theorem~\ref{thm:two-phase-expected}. The distribution is independent
of the input. The shift has density $1/(u\ln\beta)$ on $[1,\beta)$.
An entry $(q,r;\mathsf L)$ selects that initial-budget variant and
one first-batch saving minorant. Here $r=0$ means budgets
$q,k_2,k_3,\ldots$, and $r>0$ means
$q,r,k_2,k_3,\ldots$. All entries use
\[
 \mathsf A_q=M_q-\delta wq,\qquad
 \mathsf C_q=c_*M_q-d_*\delta wq.
\]
A mixture $\theta(q,r;\mathsf L)+(1-\theta)(q',r';\mathsf L')$
is conditional on that row's shift range. It compares two actual
variants; it never adds two savings from the same schedule.

Read the certificate in two stages: the small table supplies the
conditional mixing weights, and the row tables select which variants
to compare on each shift range. A row with one option assigns it
all conditional mass.
The twelve rational mixing weights are listed as integer numerators
with common denominator $10^9$:
\begin{center}
\small
\begin{tabular}{r|rrrrrr}
$i$ &1&2&3&4&5&6\\\hline
$10^9\theta_i$ &
363738471&165441165&467444028&619060535&88721951&689586436\\
$i$ &7&8&9&10&11&12\\\hline
$10^9\theta_i$ &751593847&874972087&599476351&160123414&234739978&835733626\\
\end{tabular}
\end{center}
The shift ranges also split at the integer breakpoints of $k_2$.
Thus each $k_j=\lfloor\beta^j u\rfloor$, $j=0,1,2$, is constant
inside each listed range, and the signed correction of
\eqref{eq:two-phase-correction} is constant there.

\begin{center}\small
\begin{tabular}{rll}
\hline Row&Shift range&Variant and saving minorant\\\hline
1&$[1,13/\beta^2)$&$(3,0;\mathsf C_{3})$\\
2&$[13/\beta^2,14/\beta^2)$&$\theta_{1}(3,0;\mathsf C_{3})+(1-\theta_{1})(4,0;\mathsf A_{4})$\\
3&$[14/\beta^2,4/\beta)$&$(4,0;\mathsf A_{4})$\\
4&$[4/\beta,15/\beta^2)$&$(4,0;\mathsf A_{4})$\\
5&$[15/\beta^2,16/\beta^2)$&$(4,0;\mathsf A_{4})$\\
6&$[16/\beta^2,17/\beta^2)$&$\theta_{2}(1,5;\mathsf C_{1})+(1-\theta_{2})(4,0;\mathsf A_{4})$\\
7&$[17/\beta^2,5/\beta)$&$(1,5;\mathsf C_{1})$\\
8&$[5/\beta,18/\beta^2)$&$(1,5;\mathsf C_{1})$\\
9&$[18/\beta^2,19/\beta^2)$&$(1,5;\mathsf C_{1})$\\
10&$[19/\beta^2,20/\beta^2)$&$(1,5;\mathsf C_{1})$\\
11&$[20/\beta^2,21/\beta^2)$&$\theta_{3}(1,5;\mathsf C_{1})+(1-\theta_{3})(1,6;\mathsf C_{1})$\\
12&$[21/\beta^2,6/\beta)$&$(1,6;\mathsf C_{1})$\\
13&$[6/\beta,22/\beta^2)$&$(1,6;\mathsf C_{1})$\\
14&$[22/\beta^2,23/\beta^2)$&$(1,6;\mathsf C_{1})$\\
15&$[23/\beta^2,24/\beta^2)$&$(1,6;\mathsf C_{1})$\\
16&$[24/\beta^2,7/\beta)$&$(1,7;\mathsf C_{1})$\\
17&$[7/\beta,2)$&$\theta_{4}(1,6;\mathsf C_{1})+(1-\theta_{4})(1,7;\mathsf C_{1})$\\
18&$[2,25/\beta^2)$&$\theta_{5}(1,6;\mathsf C_{1})+(1-\theta_{5})(2,7;\mathsf C_{2})$\\
19&$[25/\beta^2,26/\beta^2)$&$(2,7;\mathsf C_{2})$\\
20&$[26/\beta^2,27/\beta^2)$&$(2,7;\mathsf C_{2})$\\
21&$[27/\beta^2,28/\beta^2)$&$\theta_{6}(2,7;\mathsf C_{2})+(1-\theta_{6})(2,8;\mathsf C_{2})$\\
\hline\end{tabular}\end{center}
\begin{center}\small
\begin{tabular}{rll}
\hline Row&Shift range&Variant and saving minorant\\\hline
22&$[28/\beta^2,8/\beta)$&$(2,8;\mathsf C_{2})$\\
23&$[8/\beta,29/\beta^2)$&$(2,8;\mathsf C_{2})$\\
24&$[29/\beta^2,30/\beta^2)$&$(2,8;\mathsf C_{2})$\\
25&$[30/\beta^2,31/\beta^2)$&$(2,8;\mathsf C_{2})$\\
26&$[31/\beta^2,9/\beta)$&$(2,9;\mathsf C_{2})$\\
27&$[9/\beta,32/\beta^2)$&$\theta_{7}(2,8;\mathsf C_{2})+(1-\theta_{7})(2,9;\mathsf C_{2})$\\
28&$[32/\beta^2,33/\beta^2)$&$(2,9;\mathsf C_{2})$\\
29&$[33/\beta^2,34/\beta^2)$&$(2,9;\mathsf C_{2})$\\
30&$[34/\beta^2,35/\beta^2)$&$\theta_{8}(2,9;\mathsf C_{2})+(1-\theta_{8})(2,10;\mathsf C_{2})$\\
31&$[35/\beta^2,10/\beta)$&$(2,10;\mathsf C_{2})$\\
32&$[10/\beta,36/\beta^2)$&$\theta_{9}(2,10;\mathsf C_{2})+(1-\theta_{9})(3,10;\mathsf C_{3})$\\
33&$[36/\beta^2,37/\beta^2)$&$(3,10;\mathsf C_{3})$\\
34&$[37/\beta^2,3)$&$(3,10;\mathsf C_{3})$\\
35&$[3,38/\beta^2)$&$(3,10;\mathsf C_{3})$\\
36&$[38/\beta^2,11/\beta)$&$\theta_{10}(3,10;\mathsf C_{3})+(1-\theta_{10})(3,11;\mathsf C_{3})$\\
37&$[11/\beta,39/\beta^2)$&$(3,10;\mathsf C_{3})$\\
38&$[39/\beta^2,40/\beta^2)$&$(3,11;\mathsf C_{3})$\\
39&$[40/\beta^2,41/\beta^2)$&$(3,11;\mathsf C_{3})$\\
40&$[41/\beta^2,42/\beta^2)$&$(3,11;\mathsf C_{3})$\\
41&$[42/\beta^2,12/\beta)$&$\theta_{11}(3,11;\mathsf C_{3})+(1-\theta_{11})(3,12;\mathsf C_{3})$\\
42&$[12/\beta,43/\beta^2)$&$(3,12;\mathsf C_{3})$\\
43&$[43/\beta^2,\beta)$&$\theta_{12}(3,12;\mathsf C_{3})+(1-\theta_{12})(3,13;\mathsf C_{3})$\\
\hline\end{tabular}\end{center}

To use the table, first integrate the log-uniform shift density over
each row, then multiply by that row's conditional option weights.
The budgets are constant within a row, so its signed correction
and first-batch minorant contribute fixed coefficients. Summing
these contributions gives exactly the rank credits used in the
main proof, through the coefficient definitions in
\eqref{eq:two-phase-coefficients}.
For each row $[l,h)$ and option of conditional weight $v$, compute
$P=\ln(h/l)/\ln\beta$. Starting with $a_t=e_\beta(t)$ and $\mu=0$,
add $Pv[d(u,t)+s\mathbf1_{\{t\le q\}}]$ to $a_t$ and
$Pv\rho$ to $\mu$, where the saving minorant is
$sM_q-\rho\delta w$. The main text gives the complete finite-rank,
zero-length, and tail checks. The file
\texttt{expected-two-phase-}\allowbreak\texttt{check.py}, using the same rational data in
\texttt{expected-two-phase-}\allowbreak\texttt{certificate.json}, verifies the row
coefficients and certificate inequalities from the displayed data.

\paragraph{Complete coefficient accumulation.}
Index the rows by $i$ and the options in a row by $j$. Let the
row be $[l_i,h_i)$ and its conditional option weights be $v_{ij}$.
Put
\[
 P_i=\int_{l_i}^{h_i}\frac{du}{u\ln\beta}
     =\frac{\ln(h_i/l_i)}{\ln\beta},\qquad
 \sum_iP_i=1,\quad \sum_jv_{ij}=1.
\]
For the first identity, the antiderivative is $\ln u/\ln\beta$.
Summing the consecutive row integrals then gives
$\ln(\beta/1)/\ln\beta=1$; each conditional option sum is one by
the listed complementary mixing weights. For an option with
budgets $(q_{ij},r_{ij})$, the minorant coefficients are
\[
 (s_{ij},\rho_{ij})=
 \begin{cases}
 (1,q_{ij}),&\mathsf A_{q_{ij}},\\
 (c_*,d_*q_{ij}),&\mathsf C_{q_{ij}},
 \end{cases}
 \qquad c_*=\frac{228683}{103683},\quad d_*=\frac{250000}{103683}.
\]
These rational values are obtained by substituting $p_*$ in
$c_*=(2-p_*)/(1-p_*)$ and $d_*=2/(1-p_*)$.
Write $k_{0i},k_{1i},k_{2i}$ for that row's constant geometric
budgets and choose any $u_i\in[l_i,h_i)$. The full sums are
\[
 \begin{aligned}
 a_t&=e_\beta(t)+\sum_iP_i\sum_jv_{ij}
       \left[d_{q_{ij},r_{ij}}(u_i,t)
                       +s_{ij}\mathbf1_{\{t\le q_{ij}\}}\right],\\
 \mu&=\sum_iP_i\sum_jv_{ij}\rho_{ij},\\
 D_0(k)&=\sum_iP_i\sum_jv_{ij}
       \left[q_{ij}\mathbf1_{\{q_{ij}<k\}}
             +r_{ij}\mathbf1_{\{r_{ij}<k\}}
             -k_{0i}\mathbf1_{\{k_{0i}<k\}}
             -k_{1i}\mathbf1_{\{k_{1i}<k\}}\right].
 \end{aligned}
\]
To verify that these budgets are constant throughout each half-open
row, fix $j\in\{0,1,2\}$. Since the floor is nondecreasing, its minimum
on $[l_i,h_i)$ occurs at $l_i$. Its largest value before the excluded
right endpoint is $\lceil\beta^j h_i\rceil-1$. Therefore the endpoint
identity below certifies every shift in the row:
\[
 \lfloor\beta^j l_i\rfloor
   =\lceil\beta^j h_i\rceil-1=k_{ji}.
\]

\paragraph{A worked row.}
In row 1, $1\le u<13/\beta^2$ gives
$(k_0,k_1,k_2)=(1,3,12)$, and the sole option is
$(q,r)=(3,0)$ with minorant $\mathsf C_3$.
Its probability is
\[
 P_1=\frac{\ln(13/\beta^2)}{\ln\beta}
       =\frac{\ln13}{\ln\beta}-2,
 \qquad 0.041455991537<P_1<0.041455991538.
\]
Substitution in \eqref{eq:two-phase-correction}, with these constant
budgets, cancels both terms at rank
three:
\[
 \begin{aligned}
 d_{3,0}(u,t)
 &=1-3+4\mathbf1_{\{t>1\}}
       +15\mathbf1_{\{t>3\}}-15\mathbf1_{\{t>3\}}\\
 &=-2+4\mathbf1_{\{t>1\}}.
 \end{aligned}
\]
Consequently \eqref{eq:two-phase-coefficients} and the zero-route
correction \eqref{eq:two-phase-zero-correction} give this row's
contributions as
\[
 \begin{aligned}
 \Delta a_t&=P_1\left[-2+4\mathbf1_{\{t>1\}}
                         +c_*\mathbf1_{\{t\le3\}}\right],\\
 \Delta\mu&=3d_*P_1,\\
 \Delta D_0(k)&=P_1\left[3\mathbf1_{\{3<k\}}
                 -\mathbf1_{\{1<k\}}-3\mathbf1_{\{3<k\}}\right]
                 =-P_1\mathbf1_{\{k>1\}}.
 \end{aligned}
\]
In particular, its rank-one contribution is
$P_1(c_*-2)$, its contributions at ranks two and three are
$P_1(c_*+2)$, and every later contribution is $2P_1$.
Rows with two options use the same calculation for each option
and multiply the results by its conditional weight before adding.
For example, row 2 contributes
$P_2[\theta_1(c_*-2)+(1-\theta_1)(1-3)]$ to $a_1$ and
$P_2[3d_*\theta_1+4(1-\theta_1)]$ to $\mu$.
Each execution chooses one valid alternative. Consequently, the
expected row contribution is the weighted average of the credit from
each alternative.
